% !PRINTED 78
\weberpage{78}
%<<CONT

\begin{equation}
        a_{\alpha_1}^{(\beta_1)}a_{\alpha_2}^{(\beta_2)}\cdots a_{\alpha_\nu}^{(\beta_\nu)}\,
        A_{\alpha_1,\alpha_2,\ldots,\alpha_\nu}^{\beta_1,\beta_2,\ldots,\beta_\nu}.
\tag{2}
\end{equation}

Man kann durch Umstellen von Zeilen und Colonnen, wodurch höchstens das Zeichen der Determinante geändert wird, immer erreichen, dass die Elemente
\begin{equation}
        a_{\alpha_1}^{(\beta_1)},
        a_{\alpha_2}^{(\beta_2)},\ldots,
        a_{\alpha_\nu}^{(\beta_\nu)}
\tag{3}
\end{equation}
an die Stelle der Elemente. % [sic] Weber druckt hier einen Punkt mitten im Satz (vor dem Display); quellentreu beibehalten
\[
        a_1^{(1)},a_2^{(2)},\ldots,a_\nu^{(\nu)}
\]
gelangen; dann aber lässt sich die Regel I. auf die Bestimmung von \(A_{\alpha_1,\alpha_2,\ldots,\alpha_\nu}^{\beta_1,\beta_2,\ldots,\beta_\nu}\) anwenden und es ergiebt sich:

\gesperrt{II. Man erhält (vom Vorzeichen abgesehen) \(A_{\alpha_1,\alpha_2,\ldots,\alpha_\nu}^{\beta_1,\beta_2,\ldots,\beta_\nu}\) als \((n-\nu)\)-reihige Determinante, wenn man in \(A\) alle Zeilen und Colonnen weglässt, die sich in einem der Elemente (3) schneiden, und die übrig bleibenden Zeilen und Colonnen in ihrer Reihenfolge stehen lässt.}

Für die Zeichenbestimmung aber ergiebt sich folgende Vorschrift.

Man ordne die unteren und die oberen Indices \(1,2,\ldots,n\) in der Weise:
\begin{equation}
        \alpha_1,\alpha_2,\ldots,\alpha_\nu,\alpha_{\nu+1},\ldots,\alpha_n,
\tag{4}
\end{equation}
\begin{equation}
        \beta_1,\beta_2,\ldots,\beta_\nu,\beta_{\nu+1},\ldots,\beta_n,
\tag{5}
\end{equation}
indem man \(\alpha_{\nu+1}\ldots\alpha_n\) und ebenso \(\beta_{\nu+1}\ldots\beta_n\) \gesperrt{der Grösse nach auf einander folgend} annimmt.

\gesperrt{III. Die in II. beschriebene \((n-\nu)\)-reihige Determinante erhält das positive oder negative Zeichen, je nachdem die beiden Anordnungen (4) und (5) der Ziffern \(1,2,\ldots,n\) beide zu derselben oder zu verschiedenen Arten gehören.}

Denn die Determinante ändert ihr Zeichen durch jede Vertauschung zweier unterer oder zweier oberer Indices. Um den allgemeinen Fall (2) auf den besonderen Fall (1) zurückzuführen, hat man so viele Transpositionen oberer und unterer Indices vorzunehmen, dass die Permutationen (4) und (5) beide in die ursprüngliche Anordnung \(1,2,3\ldots n\) übergehen, und ebenso viele Zeichenwechsel haben stattgefunden.
